\section{Contexto del problema}
La diversidad del mundo natural es enorme: se estima que existen más de 5,000 especies de mamíferos, 10,000 de aves y 30,000 de peces, pero hay más de 400,000 tipos de flores \parencite{kagglePetalsMetal}. Distinguir entre ellas puede requerir conocimiento especializado y, aun para una persona con experiencia, muchas flores son visualmente muy parecidas entre sí. Identificarlas de forma manual es lento y no escala cuando se tienen miles de fotografías. Por eso resulta valioso contar con modelos que reconozcan el tipo de flor de forma automática a partir de una imagen.

Este proyecto parte de la competencia \textit{Petals to the Metal}, de Kaggle, cuyo reto es construir un modelo que identifique el tipo de flor en un conjunto de imágenes. La competencia se limita a 104 clases de flores y consiste en un problema de clasificación multiclase de imágenes en el que el modelo debe captar detalles finos de forma, estructura y color para distinguir entre categorías cercanas.

Este problema tiene varias dificultades que lo hacen relevante desde un enfoque de redes neuronales. Primero, las clases son heterogéneas, ya que algunas solo tienen un subtipo de flor y otras agrupan muchas subvariedades. Además, las imágenes no fueron capturadas en el mismo entorno, lo que les da más variedad. Por último, el número de imágenes por clase tampoco es uniforme, por lo que, considerando la métrica de evaluación, es importante que el modelo funcione bien tanto en las clases con pocos ejemplos como en las que tienen muchos.
